\documentclass[9pt]{extarticle}
\usepackage[utf8]{inputenc}
\usepackage{amsmath}
\usepackage{booktabs}
\usepackage{longtable}
\usepackage[margin=0.5in,landscape]{geometry}
\usepackage{array}
\pagestyle{empty}
\setlength{\parindent}{0pt}
\setlength{\tabcolsep}{5pt}
\renewcommand{\arraystretch}{1.05}

% Body rows live in table2_global_params_body.tex and are regenerated by
% extract_full_model_analysis.R (section 11). Do not hand-edit data rows —
% edit the producer instead. This wrapper owns the landscape longtable
% preamble and the multi-line header.

\begin{document}

\small
\begin{longtable}{@{}l>{\raggedleft\arraybackslash}p{2.2cm}>{\raggedleft\arraybackslash}p{2.2cm}>{\raggedleft\arraybackslash}p{2.3cm}>{\raggedleft\arraybackslash}p{2.2cm}>{\raggedleft\arraybackslash}p{1.8cm}>{\raggedleft\arraybackslash}p{2.2cm}cc@{}}
\caption*{\textbf{Table 2.} Global parameters and spatial heterogeneity across model variants} \\
\toprule
\textbf{Model} & \textbf{RMSE} & \textbf{R}$^2$ & \textbf{$\beta_0$} & \textbf{$\beta_{\text{OIPC}}$} & \textbf{$\lambda_{\text{int}}$} & \textbf{GP scale} & \textbf{Knot} & \textbf{Knot} \\
 & \textbf{(\textperthousand)} & & \textbf{(\textperthousand)} & & \textbf{(km)} & \textbf{(km)} & \textbf{Slope SD} & \textbf{Int. SD} \\
 & & & & & & & & \textbf{(\textperthousand)} \\
\midrule
\endfirsthead

\multicolumn{9}{c}{Table 2 (continued)} \\
\toprule
\textbf{Model} & \textbf{RMSE} & \textbf{R}$^2$ & \textbf{$\beta_0$} & \textbf{$\beta_{\text{OIPC}}$} & \textbf{$\lambda_{\text{int}}$} & \textbf{GP scale} & \textbf{Knot} & \textbf{Knot} \\
 & \textbf{(\textperthousand)} & & \textbf{(\textperthousand)} & & \textbf{(km)} & \textbf{(km)} & \textbf{Slope SD} & \textbf{Int. SD} \\
 & & & & & & & & \textbf{(\textperthousand)} \\
\midrule
\endhead

\bottomrule
\endlastfoot

% Auto-generated by extract_full_model_analysis.R.
% Do not hand-edit — regenerate via `make tables`.

baseline & 21.3 [21.2, 21.3] & 0.695 [0.693, 0.696] & -177.9 [-179.2, -176.7] & 0.778 [0.748, 0.808] & 6.4 [2.6, 14.0] & - & - & -\\
baseline$\_{\text{sp}}$ & 16.5 [16.2, 16.7] & 0.817 [0.811, 0.822] & -186.3 [-194.3, -178.7] & 0.586 [0.480, 0.688] & 5.8 [2.4, 12.0] & 3920 [2387, 5944] & 0.369 & 39.4\\
baseline\_env & 20.3 [20.2, 20.4] & 0.721 [0.718, 0.723] & -178.3 [-179.5, -177.1] & 0.860 [0.824, 0.896] & 7.9 [3.5, 15.9] & - & - & -\\
baseline\_env$\_{\text{sp}}$ & 16.2 [15.9, 16.4] & 0.823 [0.817, 0.828] & -187.7 [-196.3, -179.5] & 0.570 [0.439, 0.691] & 6.3 [2.6, 13.1] & 3877 [2362, 5961] & 0.357 & 41.3\\
baseline\_veg & 20.1 [20.0, 20.2] & 0.727 [0.725, 0.729] & -167.1 [-174.6, -159.7] & 0.699 [0.596, 0.802] & 7.9 [3.0, 19.3] & - & - & -\\
baseline\_veg$\_{\text{sp}}$ & 16.3 [16.1, 16.5] & 0.821 [0.815, 0.826] & -178.2 [-187.9, -168.6] & 0.642 [0.512, 0.772] & 6.7 [2.7, 14.4] & 3538 [1985, 5807] & 0.355 & 30.3\\
full & 20.1 [20.0, 20.2] & 0.727 [0.724, 0.730] & -171.9 [-178.1, -165.8] & 0.859 [0.817, 0.900] & 8.1 [3.5, 16.5] & - & - & -\\
full$\_{\text{sp}}$ & 16.1 [15.9, 16.4] & 0.824 [0.818, 0.829] & -179.9 [-190.2, -170.0] & 0.552 [0.426, 0.671] & 6.6 [2.7, 14.3] & 3907 [2409, 5924] & 0.334 & 40.1\\
full\_interact & 19.4 [19.3, 19.6] & 0.745 [0.742, 0.747] & -165.3 [-172.8, -157.5] & 0.727 [0.616, 0.840] & 10.5 [4.4, 22.9] & - & - & -\\
full\_interact$\_{\text{sp}}$ & 16.1 [15.9, 16.4] & 0.824 [0.819, 0.829] & -178.5 [-188.9, -168.6] & 0.632 [0.477, 0.778] & 7.0 [2.8, 14.9] & 3581 [2073, 5834] & 0.356 & 31.6\\
elevation\_only$\_{\text{sp}}$ & 16.3 [16.0, 16.5] & 0.821 [0.816, 0.826] & -186.5 [-194.3, -178.7] & 0.611 [0.489, 0.726] & 6.5 [2.7, 13.4] & 3872 [2315, 5999] & 0.390 & 38.1\\
elevation\_c4$\_{\text{sp}}$ & 16.3 [16.0, 16.5] & 0.821 [0.815, 0.826] & -186.6 [-194.8, -178.5] & 0.609 [0.485, 0.726] & 6.5 [2.7, 13.7] & 3867 [2311, 5970] & 0.390 & 38.5\\
c4\_only$\_{\text{sp}}$ & 16.5 [16.2, 16.7] & 0.817 [0.811, 0.822] & -185.2 [-193.2, -177.6] & 0.584 [0.477, 0.687] & 5.8 [2.4, 11.9] & 3886 [2338, 5950] & 0.370 & 38.3\\
elevation\_c4\_interact$\_{\text{sp}}$ & 16.2 [16.0, 16.5] & 0.822 [0.816, 0.827] & -185.4 [-194.1, -177.0] & 0.587 [0.458, 0.712] & 6.6 [2.8, 13.6] & 3890 [2345, 5974] & 0.385 & 39.4\\

\end{longtable}

\vspace{2mm}
\small
\textit{Note:} Values shown as posterior mean [95\% credible interval]. RMSE and R$^2$ are summarized \emph{per draw} from the linear-predictor mean $\mu$ (no observation noise); the reported mean is $\mathrm{mean}_s\big(\mathrm{RMSE}(y,\mu^{(s)})\big)$, not the RMSE evaluated at the posterior-mean prediction $\mu$-bar. The two differ by Jensen's inequality; see \texttt{manuscript/TABLES.md} for the convention and how it relates to scalar RMSE / R$^2$ values cited in the narrative. $\beta_0$ = global intercept on the original $\delta^2$H scale (back-transformed per Stan line 478: \texttt{beta\_0 * d2H\_sd + d2H\_mean}); $\beta_{\text{OIPC}}$ = global slope with respect to $\delta^2$H$_{\text{precip}}$ (standardized scale); $\lambda_{\text{int}}$ = integration length scale (\texttt{lambda\_decay}, km); GP scale = Gaussian process length scale (\texttt{ls\_intercept\_km}). Knot SDs are posterior means of \texttt{sigma\_slope\_spatial} (unitless) and \texttt{sigma\_intercept\_spatial} (\textperthousand), quantifying spatial heterogeneity in slope and intercept parameters across the 125 predictive-process knots.

\end{document}
